% Generated draft: requires mathematical and source review.
\documentclass[11pt]{article}
\usepackage[margin=25mm]{geometry}
\usepackage{fontspec}
\setmainfont{Latin Modern Roman}
\usepackage{amsmath,amssymb,mathtools,unicode-math}
\setmathfont{Latin Modern Math}
\usepackage{enumitem,needspace,xurl,hyperref}
\setlength{\parindent}{0pt}
\setlength{\parskip}{5pt}
\setlength{\emergencystretch}{3em}
\newcommand{\N}{\mathbb{N}}
\newcommand{\Z}{\mathbb{Z}}
\newcommand{\Q}{\mathbb{Q}}
\newcommand{\R}{\mathbb{R}}
\newcommand{\C}{\mathbb{C}}
\newcommand{\F}{\mathbb{F}}
\newcommand{\A}{\mathbb{A}}
\newcommand{\PP}{\mathbb{P}}
\newcommand{\E}{\mathbb{E}}
\newcommand{\eps}{\varepsilon}
\newcommand{\dd}{\,\mathrm{d}}
\begin{document}
\section*{UnsolvedMath ID: 20000308}
\textbf{Source problem code:} catalogue entry
\subsection{Definitions and mathematical statement}
In the intended characteristic-zero setting, a normal $\Q$-Gorenstein variety $X$ is log canonical if every prime divisor $E$ over $X$ has log discrepancy $A_X(E)\ge0$, where $A_X(E)=1+\operatorname{coeff}_E(K_Y-f^*K_X)$ on a birational model $f:Y\to X$. It is klt if all these log discrepancies are positive. A reduced complex variety is Du Bois if the natural map $\mathcal O_X\to\underline\Omega_X^0$ to the degree-zero Deligne--Du Bois complex is a quasi-isomorphism. The original statement does not name the ground field; this characteristic-zero/complex framework must be part of its interpretation rather than an asserted arbitrary-characteristic theorem.

\textbf{Original question (Schwede).} If $X$ has log-canonical singularities, must $X$ have Du Bois singularities? The source directs readers to its labels [5] for historical discussion and progress, [4], introduction, for discussion, and [2], Chapter 12. The associated bibliography is not included.

\textbf{Attached research-aid account, requiring verification.} The aid attributes to Kollár and Kovács an affirmative result for characteristic-zero log-canonical pairs $(X,\Delta)$: the underlying $X$ is Du Bois, and every reduced closed union of log-canonical centers is Du Bois. It attributes extensions to quasi-log-canonical pairs and semi-log-canonical strata to Fujino and Liu, without exact statements or references. Its mentioned ``formal Du Bois-pair refinement'' is not specified in the summary and cannot be reconstructed uniquely as a new theorem here.

The aid additionally reports the following sharp worked family. For every pair of integers
\[b\ge r\ge1,\qquad b\ge2,\]
there are algebraic rational surface germs $X_{b,r}$ with resolution graph a star: the central exceptional rational curve $E_0$ has self-intersection $-b$, the $r$ exceptional rational leaves $E_1,\ldots,E_r$ have self-intersection $-2$, each leaf meets $E_0$ once, and distinct leaves do not meet. The report claims that every member is Cohen--Macaulay, Du Bois, and $\Q$-Gorenstein, and that the central log discrepancy is
\[A_{X_{b,r}}(E_0)=\frac{4-r}{2b-r}.\]
It gives the classification within this family:
\[r\le3:\ \mathrm{klt};\qquad r=4:\ \mathrm{lc\ but\ not\ klt};\qquad r\ge5:\ \mathrm{not\ lc}.\]
Here rational surface singularity means that the first higher direct image of $\mathcal O_Y$ under a resolution vanishes. This family is reported as a converse obstruction: Du Bois, Cohen--Macaulay and $\Q$-Gorenstein together do not force log canonicity when $r\ge5$. A graph and the central discrepancy alone are not a supplied algebraic construction or a verification of discrepancies on all divisors. The coefficient field, construction, and proof of the full asserted classification are absent from this record.
\subsection{Short English statement}
Schwede asks whether log-canonical singularities are Du Bois. The aid attributes the affirmative characteristic-zero theorem to Kollár and Kovács, with later extensions, and reports a rational surface family showing that the converse fails. All parameters, the central discrepancy, and the klt/lc transition are retained as unreviewed source claims.
\subsection{Sources}
\begin{itemize}
\item Local UnsolvedMath registry, permanent ID 20000308, source code catalogue entry; \texttt{aim-algebraic-geometry-notes.json}, notes[307], tag \texttt{problem}.
\item AIM workshop \emph{Numerical invariants of singularities and higher-dimensional algebraic varieties}, Problem 5 (Schwede): \url{https://aimath.org/WWN/singularvariety/problems_final.pdf}. Original reference labels [5], [4] introduction, and [2] Chapter 12 are retained, but their bibliography is not supplied. No PDF was consulted.
\item Attached machine-generated research attempt 1 names Kollár--Kovács and Fujino--Liu without publication locators. Result \texttt{partial\_result}, problem status at run \texttt{solved}, rigor label \texttt{proved}, novelty confidence low. It is not peer reviewed and requires independent verification before citation.
\end{itemize}
\end{document}
